% This is text associated with the open textbook "Open Electromagnetics"
% See immediately below for author, date
% This work is licensed under the Creative Commons Attribution-ShareAlike 4.0 International License. (CC BY-SA 4.0) 
% To view a copy of the license, visit http://creativecommons.org/licenses/by-sa/4.0/

%ORB TITLE Magnetic Energy
%ORB AUTHOR 0        # S.W. Ellingson, Virginia Tech (ellingson@vt.edu)
%ORB DATE 20180806   # yyyymmdd
%ORB ABSTRACT Magnetic Energy

\label{m0127_Magnetic_Energy} % <-- Must match name of this file and module directory name.  Don't mess with this!
\index{energy!magnetostatic|(}

Consider a structure exhibiting inductance; i.e., one that is able to store energy in a magnetic field in response to an applied current.
This structure could be a coil, or it could be one of a variety of inductive structures that are not explicitly intended to be an inductor; for example, a coaxial transmission line.
When current is applied, the current-bearing elements of the structure exert forces on each other.
Since these elements are not normally free to move, we may interpret this force as potential energy stored in the magnetic field associated with the current (Section~\ref{m0123_Inductance}).

We now want to know how much energy is stored in this field. 
The answer to this question has relevance in several engineering applications. 
One issue is that any system that includes inductance is using some fraction of the energy delivered by the power supply to energize this inductance. 
In many electronic systems -- in power systems in particular -- inductors are periodically energized and de-energized at a regular rate.  Since power is energy per unit time, this consumes power.  Therefore, energy storage in inductors contributes to the power consumption of electrical systems.

The stored energy is most easily determined using circuit theory concepts.
First, we note that the electrical potential difference $v(t)$ (units of V) across an inductor is related to the current $i(t)$ (units of A) through the inductor as follows (Section~\ref{m0123_Inductance}):
%
\begin{equation}
v(t) = L\frac{d}{dt}i(t)
\end{equation}
%
where $L$ (units of H) is the inductance.
The instantaneous power associated with the device is
%
\begin{equation}
p(t) = v(t)i(t)
\end{equation}
%
Energy (units of J) is power (units of J/s) integrated over time. 
Let $W_m$ be the energy stored in the inductor.
At some time $t_0$ in the past, $i(t_0)=0$ and $W_m=0$.
As current is applied, $W_m$ increases monotonically.
At the present time $t$, $i(t)=I$. 
Thus, the present value of the magnetic energy is:
%
\begin{equation}
W_m = \int_{t_0}^{t_0+t} p(\tau) d\tau 
\end{equation}
%
Now evaluating this integral using the relationships established above:
%
\begin{flalign}
W_m &= \int_{t_0}^{t+t_0} v(\tau) i(\tau) d\tau \nonumber \\
    &= \int_{t_0}^{t+t_0} \left[ L\frac{d}{d\tau}i(\tau) \right] i(\tau) d\tau \nonumber \\
    &= L \int_{t_0}^{t+t_0} \left[\frac{d}{d\tau}i(\tau)\right] i(\tau) d\tau 
\end{flalign}
%
Changing the variable of integration from $\tau$ (and $d\tau$) to $i$ (and $di$) we have
%
\begin{flalign}
W_m &= L \int_{t_0}^{t+t_0} \frac{di}{d\tau}~i~d\tau \nonumber \\
    &= L \int_{0}^{I} i~di
\end{flalign}
%
Evaluating the integral we obtain the desired expression
%
\begin{equation}
\boxed{
W_m = \frac{1}{2}LI^2 
}
\label{m0127_eWm}
\end{equation}
%
%ORB HIGHLIGHT
\begin{mdframed}[backgroundcolor=yellow!20,nobreak=true] 
The energy stored in an inductor in response to a steady current $I$ is Equation~\ref{m0127_eWm}.
This energy increases in proportion to inductance and in proportion to the square of current.
\end{mdframed}

The long straight coil (Section~\ref{m0124_Inductance_of_a_Straight_Coil}) is representative of a large number of practical applications, so it is useful to interpret the above findings in terms of this structure in particular.
For this structure we found
%
\begin{equation}
L = \frac{\mu N^2 A}{l} 
\end{equation} 
%
where $\mu$ is the permeability, $N$ is the number of windings, $A$ is cross-sectional area, and $l$ is length.
The magnetic field intensity inside this structure is related to $I$ by (Section~\ref{m0120_Magnetic_Field_of_a_Coil}):
%
\begin{equation}
H = \frac{NI}{l} 
\end{equation}
%
Substituting these expressions into Equation~\ref{m0127_eWm}, we obtain
%
\begin{flalign}
W_m &= \frac{1}{2}\left[\frac{\mu N^2 A}{l}\right]\left[\frac{Hl}{N}\right]^2 \nonumber \\
    &= \frac{1}{2}\mu H^2 Al 
\end{flalign}
%
Recall that the magnetic field inside a long coil is approximately uniform.
Therefore, the density of energy stored inside the coil is approximately uniform.
Noting that the product $Al$ is the volume inside the coil, we find that this energy density is $W_m/Al$; thus:
%
\begin{equation}
w_m = \frac{1}{2} \mu H^2 
\label{m0127_eED}
\end{equation}
%
which has the expected units of energy per unit volume (J/m$^3$).

The above expression provides an alternative method to compute the total magnetostatic energy in \emph{any} structure.
Within a mathematical volume ${\mathcal V}$, the total magnetostatic energy is simply the integral of the energy density over ${\mathcal V}$; i.e., 
%
\begin{equation}
W_m = \int_{\mathcal V} w_m~dv  
\end{equation}
%
This works even if the magnetic field and the permeability vary with position.
Substituting Equation~\ref{m0127_eED} we obtain:
%
\begin{equation}
\boxed{
W_m = \frac{1}{2} \int_{\mathcal V} \mu H^2 dv  
}
\label{m0127_eEDV}
\end{equation}
%
Summarizing:
%
%ORB HIGHLIGHT
\begin{mdframed}[backgroundcolor=yellow!20,nobreak=true] 
The energy stored by the magnetic field present within any defined volume is given by Equation~\ref{m0127_eEDV}.
\end{mdframed}
%
It's worth noting that this energy increases with the permeability of the medium, which makes sense since inductance is proportional to permeability.

Finally, we reiterate that although we arrived at this result using the example of the long straight coil, Equations~\ref{m0127_eED} and \ref{m0127_eEDV} are completely general.

\index{energy!magnetostatic|)}
